\chapter{Phần mở đầu}
\label{Chapter1}

\section{Tính cấp thiết của đề tài}
\label{sec:research_urgency}
Trong khoảng một thập kỷ trở lại đây, \acrfull{ml} nói chung và \acrfull{dl} nói riêng đã đạt được những thành tựu vượt bậc trên hầu hết các bài toán trong các lĩnh vực khác nhau. Tuy nhiên, phần lớn thành công này dựa trên một giả định rằng mô hình có thể được huấn luyện trên một tập dữ liệu có kích thước lớn và đa dạng cho tác vụ cụ thể, và nhìn chung, mô hình càng sâu thì càng khai thác tốt lượng dữ liệu đó để giải quyết tác vụ. Nhưng giả định này trong bối cảnh thực tế đã vấp phải một trở ngại lớn đối với nhiều tác vụ nhạy cảm hoặc đặc thù, việc thu thập một lượng dữ liệu đủ lớn và đa dạng là rất khó khăn, thậm chí bất khả thi, do chi phí, tính đặc thù chuyên môn (y khoa, công nghiệp) hoặc các ràng buộc đạo đức và pháp lý \cite{wang2020generalizing, fawakherji2026learning}.

Bài toán \acrfull{fsl} vốn đã được đặt nền móng từ trước khi \acrlong{dl} bùng nổ, nhưng trở nên đặc biệt cấp thiết trở lại trong kỷ nguyên deep learning, khi các mô hình sâu hiện đại lại càng phụ thuộc nặng nề vào dữ liệu lớn. Các phương pháp \acrshort{fsl} dựa trên deep learning ra đời nhằm giải quyết trực tiếp trở ngại này, cho phép mô hình học và tổng quát hóa tốt chỉ với một lượng rất nhỏ mẫu dữ liệu gán nhãn. Trong số các hướng tiếp cận của bài toán \acrlong{fsl}, Meta-Learning (học cách học) vốn đã được đề xuất từ lâu, trước cả deep learning hiện đại \cite{schmidhuber1987evolutionary, bengio2013optimization}, là hướng phổ biến và đặc trưng cho lớp bài toán này. Đồng thời, với đột phá của C.Finn và các cộng sự \cite{finn2017model}, lớp phương pháp dựa trên gradient (gradient-based), đại diện tiêu biểu là \acrfull{maml} đã trở thành một trong những hướng tiếp cận chủ đạo nhờ tính linh hoạt và khả năng áp dụng cho mọi kiến trúc khả vi.

Tuy nhiên, cũng như phần lớn các mô hình học sâu khác, các mô hình meta-learning nói chung và \acrshort{maml} nói riêng vẫn là các mô hình hộp đen (blackbox models). Cơ chế mà tri thức tiền nghiệm (prior knowledge) ẩn chứa trong bộ tham số khởi tạo $\theta_0$ chi phối quá trình thích nghi nhanh (fast-adaptation) trên tập hỗ trợ (support set) $S$ vẫn chưa được biểu diễn một cách tường minh. Sự thiếu minh bạch này đặt ra rào cản lớn khi triển khai các hệ thống đòi hỏi độ tin cậy cao, nơi người dùng cần kiểm chứng được quyết định của mô hình thay vì chấp nhận nó như một phép màu thống kê \cite{lum2016predict, angwin2016machine}.

Đồng thời, tri thức tiền nghiệm cũng không phải lúc nào cũng có lợi cho tác vụ đích. Khi tác vụ đích là một tác vụ không đồng nhất (Heterogeneous task), tập hỗ trợ của tác vụ đích khó (hard), nhiễu (noisy), hoặc nằm ngoài phân phối (out-of-distribution) so với phân phối các tác vụ tại giai đoạn meta-training \cite{si2025meta}, quá trình thích nghi nhanh có thể trở nên kém hiệu quả, thậm chí làm giảm hiệu suất so với trước khi thích nghi. Mà nguyên nhân chính dẫn tới việc giảm hiệu suất chính là do sự chuyển giao tiêu cực (negative transfer) của các tri thức tiền nghiệm (prior knowledge), vốn đã được Yue và các cộng sự chứng minh trong công trình của mình \cite{yue2020interventional}.

Nhưng hầu hết các phương pháp giải quyết hiện tượng chuyển giao tiêu cực hiện chỉ được \textbf{phát hiện thông qua triệu chứng} từ các độ đo về độ chính xác giảm trên tác vụ đích, chứ chưa được \textbf{giải thích ở cấp độ nguyên nhân}, do bản thân bước thích nghi vẫn là một hộp đen. Một kết quả khác do Raghu và các cộng sự \cite{raghu2019rapid} thực hiện việc đóng băng phần thân mạng (backbone) và chỉ thực hiện thích nghi nhanh với phần đầu ra (head) của mô hình, cho thấy phần lớn hiệu quả của \acrshort{maml} đến từ việc tái sử dụng đặc trưng (feature reuse) của phần thân mô hình (backbone) đã được học từ quá trình meta-training. Điều này gợi ý rằng tập hỗ trợ $S$ đóng một vai trò lớn trong việc định hướng tín hiệu thích nghi (steering signal) trên không gian đặc trưng sẵn có. Do đó, dẫn tới một góc nhìn khác là nguyên nhân của negative transfer nhiều khả năng nằm ở việc $S$ định hướng sai vùng chú ý của mô hình trong quá trình thích nghi nhanh. Hướng lập luận này được củng cố gián tiếp bởi HSFM \cite{parast2026hsfm}, phương pháp can thiệp trực tiếp lên các vector nhúng của tập hỗ trợ $S$ (support embeddings) trong không gian đặc trưng đã được đóng băng để chống tương quan giả. 

Tuy nhiên tất cả các phương pháp trên, chỉ dừng ở việc \textbf{sửa chữa} biểu diễn nhằm tăng cường hiệu suất khái quát hóa cho tác vụ đích, mà không nhằm \textbf{giải thích cho con người} cơ chế mà $S$ đã định hướng sai ở đâu. 

Ở một nhánh khác trong lĩnh vực nghiên cứu về \acrlong{ml}, \acrfull{xai} trong những năm gần đây đã được giới học thuật chú ý trở lại và đạt được nhiều thành công đáng kể trong việc mở hộp đen cho các mô hình học sâu đơn lẻ, tiêu biểu là các phương pháp diễn giải cục bộ (local explanation) như Grad-CAM \cite{selvaraju2017grad} hay SHAP \cite{lundberg2017unified}. Tuy nhiên, việc áp dụng các kỹ thuật này vào các phương pháp học có cấu trúc tối ưu hai tầng như meta-learning/\acrshort{maml} vẫn còn rất hạn chế. Các nỗ lực hiện có theo hướng này đối với dữ liệu phi cấu trúc nói chung và dạng hình ảnh nói riêng, chủ yếu chỉ diễn giải được bước suy luận cuối cùng trên mô hình đã thích nghi (Grad-CAM cục bộ \cite{uddin2025expert}), mà bỏ qua vai trò định hướng của tập hỗ trợ $S$. Và phương pháp gần đây nổi bậc do Mitsuka và các cộng sự đề xuất \cite{mitsuka2025tlxml} nhằm mở rộng hàm ảnh hưởng (influence functions) sang cấp độ tập hợp các tác vụ (task-level) để định lượng ảnh hưởng của các tác vụ trong quá trình meta-training ảnh hưởng lên hành vi thích nghi. Nhưng nhược điểm của phương pháp này là vẫn phụ thuộc vào dữ liệu các tác vụ trong quá trình huấn luyện để tính giá trị ảnh hưởng, gây cản trở lớn về mặt chi phí tính toán, đồng thời giới hạn khả năng áp dụng trong các ứng dụng thực tế nơi dữ liệu meta-training bị đóng gói hoặc không thể truy xuất do rào cản bảo mật.

Xuất phát từ khoảng trống trên, nghiên cứu này đặt ra câu hỏi trung tâm, \textit{điều gì gây ra negative transfer trong quá trình thích nghi của meta-learning?}. Câu hỏi này có thể được cụ thể hóa thành hai câu hỏi nghiên cứu theo thứ từ lần lượt. Thứ nhất, quá trình thích nghi từ tri thức tiền nghiệm (prior knowledge) sang tri thức cụ thể (specific knowledge) với một tác vụ đích cho trước là có lợi hay có hại, và mức độ ảnh hưởng đó là bao nhiêu? Thứ hai, những đặc trưng nào trong tập hỗ trợ $S$ là nguyên nhân gây ra ảnh hưởng tích cực hay tiêu cực đó, và chúng nằm ở đâu trong không gian đặc trưng $F$?

Để giải quyết và trả lời hai câu hỏi nghiên cứu này, không chỉ đòi hỏi một độ đo định lượng mà còn yêu cầu một cơ chế diễn giải hậu nghiệm nhằm truy vết ngược từ độ đo đó về từng đặc trưng đầu vào. 

\section{Mục tiêu nghiên cứu}
\label{sec:objective}
Xuất phát từ khoảng trống đã được phân tích ở phần \ref{sec:research_urgency} với hai câu hỏi nghiên cứu đã nêu, mục tiêu của đề tài hướng đến việc xây dựng một khung diễn giải hậu nghiệm (post-hoc explanation framework) gồm hai thành phần chính, cùng với các thực nghiệm chứng minh tính hiệu quả và độ tin cậy của khung phương pháp này.

Với mục tiêu thứ nhất, là xây dựng một độ đo định lượng ứng với thành phần đầu tiên trong khung diễn giải, nhằm xác định mức độ lợi/hại của bước thích nghi nhanh trên một tác vụ đích cụ thể. Đồng thời, độ đo này phải nhạy cảm với sự thay đổi của $S$, thay vì chỉ đo hiệu suất tổng thể như độ chính xác truyền thống.

Với mục tiêu thứ hai, là xây dựng một phương pháp diễn giải thông qua trực quan hóa nhằm định vị các đặc trưng trong $S$, ứng với vị trí trong không gian đặc trưng chung $F$, đóng góp tích cực hay tiêu cực cho độ lợi thích nghi đo được ở thành phần thứ nhất.

Cuối cùng, đề tài đặt mục tiêu là thông qua các thực nghiệm định lượng lẫn định tính trên các bộ dữ liệu và backbone chuẩn, nhằm minh chứng phương pháp được đề xuất là hiệu quả, chính xác (correctness) và đáng tin cậy (trustworthiness) trong việc phản ánh đúng cơ chế thích nghi.

\section{Phương pháp tiếp cận}
\label{sec:proposed_method}

Tương ứng với các mục tiêu đã nêu, nghiên cứu đề xuất hai giải pháp kế tiếp nhau, cùng hai phép kiểm định đảm bảo tính đúng đắn và đáng tin cậy của cả hai. Đối với mục tiêu thứ nhất, độ lợi thích nghi (adaptation gain) $\Delta M$ được xây dựng dựa trên chênh lệch giá trị mất mát trên tập truy vấn $Q$ tại thời điểm trước và sau khi thực hiện quá trình thích nghi nhanh, đóng vai trò vừa là câu trả lời định lượng câu hỏi nghiên cứu thứ nhất, vừa là hàm mục tiêu để truy vết nguyên nhân ở bước tiếp theo. Đối với mục tiêu thứ hai, dựa trên việc mở rộng kỹ thuật Class Activation Mapping sang nhiều bước gradient liên tiếp trong quá trình thích nghi, nhằm lan truyền ngược độ lợi thích nghi $\Delta M$ về từng đặc trưng trong tập hỗ trợ $S$ trên không gian đặc trưng chung, qua đó xác định đặc trưng nào có lợi, đặc trưng nào có hại cho quá trình thích nghi.

Cuối cùng, tính đúng đắn (correctness) của FAMA được kiểm định thông qua hai phép kiểm tra. Thứ nhất là \acrfull{biadt}, điều chỉnh từ các bài kiểm tra xóa-và-đo (deletion test) truyền thống sang bối cảnh động của quá trình thích nghi, nhằm kiểm chứng tính trung thực (faithfulness) của FAMA đối với hành vi thực của mô hình. Thứ hai là bài kiểm tra độ nhạy (sanity check) \cite{adebayo2018sanity}, được điều chỉnh để xác minh rằng các attribution do FAMA sinh ra thực sự phụ thuộc vào tham số đã học của mô hình, chứ không phải một phép biến đổi không mang thông tin của đầu vào.

\section{Đóng góp của đề tài}
\label{sec:contributions}

Trong khuôn khổ luận văn, đề tài mang lại các đóng góp chính sau. 

\textbf{Thứ nhất,} đề tài phát biểu negative transfer trong meta-learning như một bài toán diễn giải hậu nghiệm ở cấp độ đặc trưng cho \textit{bản thân quá trình thích nghi nhanh (fast-adaptation)}, thay vì chỉ là một triệu chứng đo bằng độ chính xác. Góc nhìn này khác với TLXML \cite{mitsuka2025tlxml} vốn vẫn coi bước thích nghi là một hộp đen và Grad-CAM cục bộ vốn chỉ diễn giải được bước suy luận cuối cùng trên mô hình đã thích nghi (Chương~\ref{Chapter3}).

\textbf{Thứ hai,} đề tài đề xuất độ lợi thích nghi $\Delta M$, là độ đo định lượng mức độ lợi/hại mà bước thích nghi mang lại trên một tác vụ đích cụ thể, làm cơ sở để phát hiện negative transfer một cách định lượng thay vì chỉ định tính (Chương~\ref{Chapter3}).

\textbf{Thứ ba,} đề tài đề xuất \acrfull{fama}, phương pháp trực quan hóa trực tiếp đặc trưng nào trong $S$ có lợi hay có hại cho quá trình thích nghi trên không gian đặc trưng chung, khác với HSFM \cite{parast2026hsfm} ở chỗ hướng đến diễn giải cho con người thay vì chỉ sửa chữa biểu diễn để tăng hiệu năng (Chương~\ref{Chapter3}).

\textbf{Thứ tư,} đề tài đề xuất \acrshort{biadt} và điều chỉnh phép sanity check sang bối cảnh thích nghi động. Hai công cụ này không phụ thuộc riêng vào \acrshort{fama} và có thể áp dụng để kiểm định bất kỳ phương pháp XAI nào khác được đề xuất trong tương lai cho bài toán thích nghi trong meta-learning (Chương~\ref{Chapter4}).

\section{Bố cục luận văn}
\label{sec:structure}
Phần còn lại của luận văn được tổ chức thành các chương như sau:

\begin{itemize}
    \item \textbf{Chương 2 - Tổng quan nghiên cứu}: trình bày nền tảng lý thuyết về \acrlong{fsl}, Meta-Learning, cơ chế thích nghi hai tầng của MAML, các phương pháp diễn giải hậu nghiệm dựa trên saliency map, và các nguyên lý kiểm định XAI (completeness, sanity check). Trên cơ sở đó, các phân tích và đánh giá của các công trình liên quan trực tiếp về cơ chế nội tại của MAML, và về XAI ở mức tác vụ (TLXML) hay feature-space (HSFM) sẽ được trình bày nhằm nêu rõ những vấn đề mà các hướng này chưa giải quyết được, từ đó chỉ ra những vấn đề mà đề tài tập trung giải quyết thông qua $\Delta M$ và FAMA.

    \item \textbf{Chương 3 - Phương pháp đề xuất}: phát biểu bài toán và trình bày chi tiết hai thành phần chính của đề tài. Gồm định nghĩa độ lợi thích nghi $\Delta M$ và phương pháp \acrfull{FAMA}, cơ sở toán học và ý nghĩa của chúng.

    \item \textbf{Chương 4 - Thực nghiệm và đánh giá}: mô tả thiết lập thực nghiệm trên các bộ dữ liệu chuẩn và các phương pháp đối chứng. Đồng thời, cũng trình bày các bài kiểm định tính đúng đắn và độ tin cậy của \acrshort{fama}, và phân tích kết quả định lượng lẫn định tính.

    \item \textbf{Chương 5 - Kết luận và hướng phát triển}: tổng kết các kết quả đạt được, hạn chế của phương pháp, và đề xuất hướng phát triển trong tương lai (mở rộng sang các phương pháp meta-learning khác ngoài MAML, sang các bài toán ngoài phân loại ảnh, và ứng dụng kết quả diễn giải vào việc giảm thiểu negative transfer trong quá trình huấn luyện).
\end{itemize}

%Ngôn ngữ để viết và trình bày báo cáo khóa luận tốt nghiệp, đồ án tốt nghiệp, thực tập tốt nghiệp (sau đây gọi chung là báo cáo) là tiếng Việt hoặc tiếng Anh. 
%Trường hợp chọn ngôn ngữ tiếng Anh để viết và trình bày báo cáo,  sinh viên cần có đơn đề nghị, được cán bộ hướng dẫn (CBHD) đồng ý và nộp cho bộ phận Giáo vụ của Khoa vào thời điểm đăng ký đề tài để xin ý kiến.
%Báo cáo viết và trình bày bằng tiếng Anh phải có bản tóm tắt viết bằng tiếng Việt.


%Tóm tắt luận văn được trình bày nhiều nhất trong 24 trang in trên hai mặt giấy, cỡ chữ Times New Roman 11 của hệ soạn thảo Winword hoặc phần mềm soạn thảo Latex đối với các chuyên ngành thuộc ngành Toán.

%Mật độ chữ bình thường, không được nén hoặc kéo dãn khoảng cách giữa các chữ.
%Chế độ dãn dòng là Exactly 17pt.
%Lề trên, lề dưới, lề trái, lề phải đều là 1.5 cm.
%Các bảng biểu trình bày theo chiều ngang khổ giấy thì đầu bảng là lề trái của trang.
%Tóm tắt luận án phải phản ảnh trung thực kết cấu, bố cục và nội dung của luận án, phải ghi đầy đủ toàn văn kết luận của luận án.
%Mẫu trình bày trang bìa của tóm tắt luận văn (phụ lục 1).
